\section{灵敏度分析与模型检验}

\subsection{灵敏度分析方法}

为评估模型关键参数对优化结果的影响程度，选取购电电价、光伏装机容量、风电装机容量和售电电价四个参数进行灵敏度分析。以问题三典型场景下54吨/日产量的最优方案为基准（基准吨氨成本4530.88元/吨），分别对各参数进行$\pm$20\%--30\%的扰动，观察吨氨成本的变化幅度。

\subsection{灵敏度分析结果}
\begin{table}[H]
\centering
\caption{关键参数灵敏度分析结果}\label{tab:sensitivity}
\renewcommand{\arraystretch}{0.8}  % 添加这行
\begin{tabular}{lcccc}
\toprule
参数 & 扰动范围 & 成本降低(元/吨) & 成本升高(元/吨) & 敏感度排名 \\
\midrule
购电电价 & $\pm$30\% & $-$501 & +501 & 1（最敏感） \\
光伏装机容量 & $\pm$20\% & $-$382 & +493 & 2 \\
风电装机容量 & $\pm$20\% & $-$259 & +266 & 3 \\
售电电价 & $\pm$30\% & $-$85 & +85 & 4（最不敏感） \\
\bottomrule
\end{tabular}
\end{table}

由表\ref{tab:sensitivity}可知，关键参数对吨氨成本的影响程度排序为：购电电价 > 光伏容量 > 风电容量 > 售电电价。

购电电价作为最敏感的参数±30\%的扰动引起可成本变动±501元/吨（±11.1\%），主因是风光不足时段的购电成本占比较高，表明低电价区域选址对项目经济性至关重要。光伏容量敏感度居次且具非对称性，源于白天出力集中易饱和，增容边际效益有限，而减容则会显著增加高价购电需求。风电容量敏感度位列第三且相对对称，主要得益于其全天出力均衡，对购电需求的影响呈近似线性。

\subsection{模型检验}

\subsubsection{能量守恒验证}

对所有计算场景进行能量守恒检验：$E_{RE} + E_{buy} = E_{total} + E_{sell}$。所有场景的能量平衡误差均小于0.001 MWh（相对误差$<$0.0002\%），验证了功率平衡模型的正确性。

\subsubsection{约束满足验证}

检验所有优化结果是否满足物理约束：（1）所有功率变量非负；（2）$\alpha(t) \in [0.1, 1.0]$（问题三、四）；（3）购售电互斥（同一时段$P_{buy}(t) \cdot P_{sell}(t) = 0$）；（4）储能SOC在$[0, C_{sto}]$范围内；（5）日产量不超过72吨/日。所有约束均严格满足。